\documentclass[12pt,a4paper]{article}
\usepackage[utf8]{inputenc}
\usepackage[spanish]{babel}
\usepackage{amsmath,amssymb,amsfonts}
\usepackage{geometry}
\geometry{margin=2.5cm}
\usepackage{enumitem}
\usepackage{hyperref}
\usepackage{titlesec}

\title{\textbf{Cuestionario de Gravitación en Medios Continuos}}
\author{}
\date{}

\begin{document}
\maketitle
\thispagestyle{empty}

\begin{enumerate}[label=\textbf{\arabic*.}, itemsep=1.5em]

% Pregunta 1
\item \textbf{¿Cuál es la ecuación fundamental que relaciona el potencial gravitacional $\Phi$ con la densidad de masa $\rho$ en un medio continuo autogravitante?} \\
\textit{Pista: Recuerde la ley de Gauss en su forma diferencial aplicada al campo gravitatorio.}
\begin{enumerate}[label=\alph*), leftmargin=*]
    \item $\nabla \cdot \vec{v} = 0$
    \item $\nabla^2 \Phi = 4\pi G \rho$
    \item $\nabla p = \rho \vec{g}$
    \item $\frac{D\vec{v}}{Dt} = 0$
\end{enumerate}

% Pregunta 2
\item \textbf{En el modelo de una losa isotérmica autogravitante infinita en el plano $xy$, ¿de qué coordenada(s) depende la densidad $\rho$ por simetría?} \\
\textit{Pista: Considere qué direcciones son indistinguibles en una geometría de plano infinito.}
\begin{enumerate}[label=\alph*), leftmargin=*]
    \item De las coordenadas polares $r$ y $\theta$
    \item De $x, y, z$ simultáneamente
    \item Únicamente de $z$
    \item Únicamente de la distancia radial al origen
\end{enumerate}

% Pregunta 3
\item \textbf{Para un gas ideal isotérmico, ¿cuál es la relación correcta entre la presión $p$ y la densidad $\rho$?} \\
\textit{Pista: Piense en la ley de los gases ideales cuando la temperatura se mantiene constante.}
\begin{enumerate}[label=\alph*), leftmargin=*]
    \item $p = \rho^\gamma$
    \item $\nabla p = 0$
    \item $p = \frac{1}{2} \rho v^2$
    \item $p = c_s^2 \rho$
\end{enumerate}

% Pregunta 4
\item \textbf{¿Qué condición física define el equilibrio hidrostático en la Clase 6?} \\
\textit{Pista: Considere las fuerzas que actúan sobre un elemento de masa que no se mueve.}
\begin{enumerate}[label=\alph*), leftmargin=*]
    \item El balance entre la fuerza del gradiente de presión y la fuerza gravitacional
    \item La conservación de la energía cinética del fluido
    \item Que la divergencia de la velocidad sea siempre positiva
    \item Que la viscosidad anule los gradientes de densidad
\end{enumerate}

% Pregunta 5
\item \textbf{Al estudiar una atmósfera planetaria sin autogravedad, ¿qué suposición se hace sobre el campo gravitatorio $\vec{g}$?} \\
\textit{Pista: Piense en cómo tratamos la gravedad terrestre al analizar fenómenos cerca de la superficie.}
\begin{enumerate}[label=\alph*), leftmargin=*]
    \item Se calcula mediante la ecuación de Poisson usando la densidad de la atmósfera
    \item Se considera un campo externo constante
    \item Se asume que $\vec{g} = 0$ en toda la atmósfera
    \item Se considera que varía linealmente con la densidad local
\end{enumerate}

% Pregunta 6
\item \textbf{En el modelo de losa isotérmica autogravitante, la ecuación combinada de equilibrio y Poisson conduce a una solución para la densidad de la forma:} \\
\textit{Pista: Busque la función que describe una estructura achatada con simetría respecto al plano central.}
\begin{enumerate}[label=\alph*), leftmargin=*]
    \item $\rho(z) = \rho_0 e^{-z/H}$
    \item $\rho(z) = \rho_0 \left(1 - \frac{z^2}{R^2}\right)$
    \item $\rho(z) = \rho_0 \text{sech}^2\left(\frac{z}{H}\right)$
    \item $\rho(z) = \text{constante}$
\end{enumerate}

% Pregunta 7
\item \textbf{Si se tiene una línea de masa de longitud $2a$ y densidad lineal $\lambda$, ¿cómo se comporta el potencial gravitacional $\Phi$ a grandes distancias ($r \gg a$)?} \\
\textit{Pista: Considere el principio de aproximación de multipolos para objetos distantes.}
\begin{enumerate}[label=\alph*), leftmargin=*]
    \item De forma logarítmica respecto a la distancia
    \item Se vuelve constante e igual a cero en todo el espacio
    \item Aumenta linealmente con la distancia $r$
    \item Como el potencial de una masa puntual $M = 2a\lambda$
\end{enumerate}

% Pregunta 8
\item \textbf{¿Qué sucede con el campo gravitatorio de una línea de masa cuando la longitud $a$ tiende a infinito?} \\
\textit{Pista: Piense en la analogía con el campo eléctrico de un hilo infinito cargado.}
\begin{enumerate}[label=\alph*), leftmargin=*]
    \item El campo se vuelve puramente radial (en coordenadas cilíndricas) y escala como $1/r$
    \item El campo desaparece debido a la dilución de la masa
    \item El campo se vuelve uniforme y constante en todas las direcciones
    \item El potencial diverge positivamente en el infinito
\end{enumerate}

% Pregunta 9
\item \textbf{¿Cuál es la expresión para la aceleración de la gravedad $g(z)$ dentro de una losa plana infinita de densidad constante $\rho_0$?} \\
\textit{Pista: Aplique la ley de Gauss usando una 'caja' que cruce el plano central.}
\begin{enumerate}[label=\alph*), leftmargin=*]
    \item $g(z) = -g_0$ (constante)
    \item $g(z) = -4\pi G \rho_0 z$
    \item $g(z) = -\frac{GM}{z^2}$
    \item $g(z) = 0$
\end{enumerate}

% Pregunta 10
\item \textbf{En el contexto de la Clase 6, ¿qué representa la escala de altura $H$ en una atmósfera isotérmica curva?} \\
\textit{Pista: Piense en el argumento de la función exponencial de caída.}
\begin{enumerate}[label=\alph*), leftmargin=*]
    \item El radio total del planeta incluyendo su atmósfera
    \item La velocidad máxima de las partículas del gas
    \item La distancia vertical sobre la cual la densidad disminuye por un factor de $e$
    \item El punto donde la presión del gas se iguala a la presión del vacío
\end{enumerate}

% Pregunta 11
\item \textbf{En la derivación del perfil de una losa autogravitante, ¿por qué se utiliza la condición $\frac{d\Phi}{dz} = 0$ en $z=0$?} \\
\textit{Pista: Imagine la dirección de la fuerza de gravedad justo en el medio de la losa.}
\begin{enumerate}[label=\alph*), leftmargin=*]
    \item Porque la presión es mínima en el centro
    \item Debido a que el potencial debe ser infinito en el origen
    \item Para asegurar que el fluido sea incompresible
    \item Por la simetría del problema respecto al plano central
\end{enumerate}

% Pregunta 12
\item \textbf{Si una atmósfera planetaria es 'delgada', ¿qué simplificación es válida para el potencial gravitacional externo $\Phi_{ext}$?} \\
\textit{Pista: Piense en la expansión de Taylor de la energía potencial cerca de la superficie.}
\begin{enumerate}[label=\alph*), leftmargin=*]
    \item Se puede aproximar como $\Phi_{ext} \approx g_0 z$
    \item Debe mantenerse estrictamente como $-GM/r$
    \item Se ignora por completo
    \item Se reemplaza por una fuerza centrífuga
\end{enumerate}

% Pregunta 13
\item \textbf{¿Cuál es la relación entre el campo gravitatorio $\vec{g}$ y el potencial $\Phi$?} \\
\textit{Pista: Considere cómo se obtiene una fuerza a partir de una energía potencial.}
\begin{enumerate}[label=\alph*), leftmargin=*]
    \item $\vec{g} = \nabla \times \vec{A}$
    \item $\vec{g} = -\nabla \Phi$
    \item $\vec{g} = \Phi^2$
    \item $\vec{g} = \frac{\partial \Phi}{\partial t}$
\end{enumerate}

% Pregunta 14
\item \textbf{En un sistema autogravitante, ¿por qué la densidad no puede ser uniforme en todo el espacio infinito?} \\
\textit{Pista: Piense en lo que sucedería con la fuerza total si hubiera masa en todas partes.}
\begin{enumerate}[label=\alph*), leftmargin=*]
    \item Debido al principio de exclusión de Pauli
    \item Porque violaría la conservación de la masa
    \item Porque el potencial gravitatorio divergiría y no habría equilibrio con la presión
    \item Porque el gas ideal siempre debe expandirse
\end{enumerate}

% Pregunta 15
\item \textbf{¿Cómo se define la masa por unidad de área $\Sigma$ para una losa infinita en el eje $z$?} \\
\textit{Pista: Piense en sumar todas las contribuciones de masa en una columna vertical.}
\begin{enumerate}[label=\alph*), leftmargin=*]
    \item $\Sigma = \rho \times \text{Volumen}$
    \item $\Sigma = \frac{d\rho}{dz}$
    \item $\Sigma = 2\pi r \rho$
    \item $\Sigma = \int_{-\infty}^{\infty} \rho(z) dz$
\end{enumerate}

% Pregunta 16
\item \textbf{Para el modelo de losa isotérmica, ¿qué ocurre con el ancho efectivo de la distribución si la temperatura aumenta (mayor $c_s$)?} \\
\textit{Pista: Considere la competencia entre la presión (térmica) y la gravedad.}
\begin{enumerate}[label=\alph*), leftmargin=*]
    \item La losa se vuelve más gruesa (aumenta $H$)
    \item La losa colapsa hacia un plano de grosor nulo
    \item La densidad central $\rho_0$ aumenta indefinidamente
    \item No hay ningún cambio por ser un sistema infinito
\end{enumerate}

% Pregunta 17
\item \textbf{¿Qué tipo de coordenadas son más apropiadas para resolver el potencial de una 'cuerda cósmica' o línea de masa infinita?} \\
\textit{Pista: Piense en la forma geométrica del objeto descrito.}
\begin{enumerate}[label=\alph*), leftmargin=*]
    \item Coordenadas esféricas
    \item Coordenadas cilíndricas
    \item Coordenadas cartesianas con dependencia en $x, y, z$
    \item Coordenadas intrínsecas de trayectoria
\end{enumerate}

% Pregunta 18
\item \textbf{En la Clase 6, al analizar una atmósfera curva sin autogravedad, ¿cómo varía la gravedad con la altura $z$ si el radio del planeta es $R$?} \\
\textit{Pista: Aplique la ley de Newton para la fuerza gravitatoria de un cuerpo esférico.}
\begin{enumerate}[label=\alph*), leftmargin=*]
    \item $g(z) = g_0 (1 + z/R)$
    \item $g(z) = g_0 e^{-z/R}$
    \item $g(z) = g_0 \frac{R^2}{(R+z)^2}$
    \item $g(z) = \text{constante}$
\end{enumerate}

% Pregunta 19
\item \textbf{En un medio autogravitante, ¿qué sucede si se duplica la constante de gravitación $G$ manteniendo la temperatura igual?} \\
\textit{Pista: Considere cómo afecta la fuerza de gravedad a la estructura del medio.}
\begin{enumerate}[label=\alph*), leftmargin=*]
    \item La presión central disminuye a la mitad
    \item La densidad se vuelve uniforme en todo el espacio
    \item El sistema se expande indefinidamente
    \item La distribución de masa se vuelve más compacta
\end{enumerate}

% Pregunta 20
\item \textbf{¿Cuál es la unidad de medida de la densidad de fuerza de cuerpo $f_e$ utilizada en la ecuación $\nabla p = f_e$?} \\
\textit{Pista: Recuerde que el gradiente de presión tiene unidades de fuerza sobre volumen.}
\begin{enumerate}[label=\alph*), leftmargin=*]
    \item $N/m^3$ (Fuerza por unidad de volumen)
    \item $kg/m^3$
    \item Pascal ($Pa$)
    \item $m/s^2$
\end{enumerate}

% Pregunta 21
\item \textbf{En el estudio de una losa isotérmica, la integración de $\frac{d^2\Phi}{dz^2}$ asumiendo equilibrio requiere conocer:} \\
\textit{Pista: Identifique las variables presentes en la ecuación de estado y la condición central.}
\begin{enumerate}[label=\alph*), leftmargin=*]
    \item La masa total del universo
    \item La densidad en el plano central y la velocidad del sonido
    \item La viscosidad cinemática del gas
    \item El color de las partículas del medio
\end{enumerate}

% Pregunta 22
\item \textbf{Verdadero o Falso: En una losa infinita autogravitante, el campo gravitatorio en el exterior de la losa ($|z| > \text{espesor}$) es constante.} \\
\textit{Pista: Compare este caso con el campo eléctrico generado por un plano infinito de carga.}
\begin{enumerate}[label=\alph*), leftmargin=*]
    \item Falso
    \item Depende de la densidad central
    \item Verdadero
    \item No se puede determinar sin conocer la temperatura
\end{enumerate}

% Pregunta 23
\item \textbf{¿Qué representa el término $\nabla \Phi$ en la física de medios continuos de la Clase 6?} \\
\textit{Pista: Relacione el operador nabla con la obtención de fuerzas a partir de energías escalares.}
\begin{enumerate}[label=\alph*), leftmargin=*]
    \item La tasa de flujo de masa
    \item El cambio de presión por unidad de temperatura
    \item La curvatura del espacio-tiempo
    \item El negativo de la aceleración de la gravedad
\end{enumerate}

% Pregunta 24
\item \textbf{¿Cuál es el valor del campo gravitatorio neto exactamente en el centro ($r=0$) de un planeta esféricamente simétrico?} \\
\textit{Pista: Considere la simetría de las fuerzas gravitatorias en el punto central.}
\begin{enumerate}[label=\alph*), leftmargin=*]
    \item Cero
    \item Infinito
    \item $g_0$
    \item Depende de la presión atmosférica
\end{enumerate}

% Pregunta 25
\item \textbf{¿Qué forma tiene la superficie de igual presión (isobara) en un fluido en equilibrio bajo un campo gravitatorio central esférico?} \\
\textit{Pista: Recuerde que en equilibrio, las isobaras coinciden con las equipotenciales.}
\begin{enumerate}[label=\alph*), leftmargin=*]
    \item Planos horizontales paralelos
    \item Esferas concéntricas
    \item Cilindros alineados con el eje de rotación
    \item Superficies irregulares dependientes de la temperatura
\end{enumerate}

% Pregunta 26
\item \textbf{En la Clase 6, el modelo de una atmósfera isotérmica plana sin autogravedad predice una densidad que decae:} \\
\textit{Pista: Recuerde la solución de la ecuación diferencial $y' = -ky$.}
\begin{enumerate}[label=\alph*), leftmargin=*]
    \item Linealmente con la altura
    \item Según el inverso del cuadrado de la distancia
    \item Exponencialmente con la altura
    \item Logarítmicamente
\end{enumerate}

\end{enumerate}

\newpage
\section*{Respuestas}

\begin{center}
\renewcommand{\arraystretch}{1.5}
\begin{tabular}{|c|c|c|c|c|c|c|c|c|c|c|}
\hline
\textbf{Preg.} & 1 & 2 & 3 & 4 & 5 & 6 & 7 & 8 & 9 & 10 \\ \hline
\textbf{Resp.} & B & C & D & A & B & C & D & A & B & C \\ \hline
\end{tabular}
\vspace{0.4cm}

\begin{tabular}{|c|c|c|c|c|c|c|c|c|c|c|}
\hline
\textbf{Preg.} & 11 & 12 & 13 & 14 & 15 & 16 & 17 & 18 & 19 & 20 \\ \hline
\textbf{Resp.} & D & A & B & C & D & A & B & C & D & A \\ \hline
\end{tabular}
\vspace{0.4cm}

\begin{tabular}{|c|c|c|c|c|c|c|}
\hline
\textbf{Preg.} & 21 & 22 & 23 & 24 & 25 & 26 \\ \hline
\textbf{Resp.} & B & C & D & A & B & C \\ \hline
\end{tabular}
\end{center}

\end{document}